\subsection{\textit{Sprint} 2}

Para a \textit{sprint} 2, o time havia previsto a criação de cerca de 6 telas com seu
\textit{back-end} respectivo, entre elas criação e gerenciamento de módulos, setores,
cursos, atividades e colaboradores, além da tela inicial. Além disso, haviam sido
previstas 2 atividades de \textit{back-end}: autenticação e envio de \textit{e-mail}. Quanto ao meu time,
havia sido previsto o desenvolvimento de um componente \textit{table}, além da página de
administrar os colaboradores.

O time desempenhou muito bem em relação às expectativas, visto o tanto de
problemas encontrados no desenvolvimento. Das 6 páginas propostas é possível
dizer que mais de 95\% do prometido foi entregue, faltando apenas detalhes de
padronização de componentes e estilização de páginas, ficando 2 \textit{front-ends} e 1
\textit{back-end} com previsão de mudanças. Quanto ao meu time, aprendemos o
componente \textit{table} e seu funcionamento, além de implementá-lo como componente
geral no projeto e construir a página de administração com tabela e filtros. Além
disso, criamos o modal da página de administração. Por fim, foi possível começar a
pensar na integração para a página de administração.

O principal problema encontrado pelo time vai além da programação. Desde
o começo do mês de maio de 2024, o estado do Rio Grande do Sul, incluindo
principalmente cidades que fazem fronteira com Porto Alegre, enfrentaram uma das
maiores enchentes da história, paralisando as aulas por um tempo e preocupando
muitos. O tempo ocioso junto com a preocupação não foi tanto problema para esta
\textit{sprint}, já que, na hora em que isso aconteceu, a maioria dos itens da entrega já estava
pronta; no entanto, tende a retardar um pouco a velocidade de desenvolvimento do
início da próxima \textit{sprint}. Além disso, entregas de trabalhos e provas foram uma
preocupação paralela à da \textit{sprint} que inevitavelmente atrapalharam um pouco o
desenvolvimento no começo da \textit{sprint}, empurrando a maior parte do
desenvolvimento para a última metade.

Cada vez mais, conciliar a Ages com estudo e estágio tem se provado uma
tarefa que requer planejamento adiantado, principalmente em uma \textit{sprint} que teve
tantas incertezas como esta, a habilidade de retraçar o planejamento e se adaptar à
situação é colocada em prova, ganhando experiência e aprendendo a agir com
calma mediante situações complicadas e imprevistos.

A equipe planeja terminar os últimos ajustes que ficaram faltando nas páginas
pendentes da \textit{sprint} passada, além de continuar desenvolvendo as telas que faltam e
adiantar bem a parte da integração. Quanto à minha parte, ainda não foi designada
nenhuma tarefa de imediato, mas já tenho em mente continuar a parte da integração
da página de administração, além de colaborar com o time nas telas que forem
posteriormente designadas.
